\documentclass[12pt, a4paper]{ctexart}
\usepackage[margin=2cm]{geometry}
\usepackage{libertine}

\usepackage{titlesec}
\usepackage{zhnumber}
\titleformat*{\section}{\Large\bfseries\raggedright}
\renewcommand{\thesection}{\zhnum{section}}
\renewcommand{\thesubsection}{\arabic{subsection}}

\setcounter{secnumdepth}{4}
\renewcommand{\theparagraph}{\thesubsubsection.\arabic{paragraph}}
\renewcommand{\paragraph}[1]{
    \refstepcounter{paragraph}
    {\bfseries\theparagraph\quad#1}\par
    \vspace{2pt}
    \noindent
}

\usepackage{enumitem}
\setlist[enumerate]{itemsep=2pt, parsep=0pt, partopsep=0pt, topsep=2pt}
\setlist[itemize]{itemsep=2pt, parsep=0pt, partopsep=0pt, topsep=2pt}
\linespread{1.2}

\usepackage{minted}
\setminted{
    breaklines=true, escapeinside=||, fontsize=\small, frame=lines,
    mathescape=$$, numbers=none, style=bw, tabsize=4
}

\usepackage{graphicx}
\usepackage{siunitx}

\usepackage{amsmath}
\usepackage{amssymb}

\usepackage{hyperref}
\hypersetup{hidelinks}

\begin{document}

\pagestyle{plain}
\thispagestyle{empty}

\noindent
\begin{tabular*}{\textwidth}{l @{\extracolsep{\fill}} r @{\extracolsep{6pt}} l}
    \LARGE{\textbf{实验报告}} & 数据结构 & \textit{Data Structures} \\
\end{tabular*}\\\\
\begin{tabular*}{\textwidth}{l l}
    \textbf{报告标题: } & \textbf{文本的查找与替换} \\
    \textbf{学号: } & 19240212 \\
    \textbf{姓名: } & 华博文 \\
    \textbf{日期: } & 2026 年 3 月 6 日 \\
\end{tabular*}\\
\rule[2ex]{\textwidth}{2pt}

\tableofcontents
\newpage

本实验的核心目标并不仅仅是 ``做出一个可以替换文本的程序''，而是把题目要求抽象为一个可维护、可扩展、可稳定运行的命令行工具。围绕这一目标，项目将任务拆解为四个层面：第一层是算法层，要求在 ``固定字符串匹配'' 和 ``正则表达式匹配'' 之间提供清晰分工，使程序既能覆盖基础需求，也能处理复杂模式；第二层是工程层，要求通过模块解耦实现清晰职责边界，避免主函数堆积逻辑导致后续迭代困难；第三层是交互层，要求命令行语义明确、参数关系合理、错误提示可读，以便在实验演示时能快速定位误用；第四层是可靠性层，要求原地写回过程具备容错性，避免出现部分写入导致的数据破坏。

在课程设计语境下，这样的目标设定体现了 ``算法与工程并重'' 的训练导向：既要解释为什么选择 KMP 与 Regex 双引擎，也要解释为何需要颜色控制、差异显示、退出码语义、临时文件原子替换等工程细节的设计。

\section{实验环境}

\subsection{操作系统}

Ubuntu 24.04.3 LTS，Linux 6.14.0-36-generic，AMD Ryzen 7 7700 (16) @ 5.3GHz，DDR5。

\subsection{编程工具}

Visual Studio Code 1.109.5，NVIM v0.12.0-dev，cmake 3.28.3，gcc 13.3.0。

\subsection{依赖库}

CLI11（\href{https://github.com/CLIUtils/CLI11}{\textit{https://github.com/CLIUtils/CLI11}}）、rang（\href{https://github.com/agauniyal/rang}{\textit{https://github.com/agauniyal/rang}}）。

\subsection{其他工具}

\LaTeX{} (\TeX\textit{studio})。

\section{算法思想描述}

\subsection{数据生成与问题建模}

本实验的目标是实现 ``从文件 A 读取文本，对模式 $s$ 完成查找与替换，再输出到文件 B'' 的完整链路。为了验证程序在不同输入条件下都能稳定工作，测试数据不采用单一文本，而是按功能点分层构造：

\begin{itemize}
    \item 基础文本：覆盖普通英文单词、重复词和无命中行；
    \item 大小写混合文本：验证 \mintinline{text}{-i} 忽略大小写行为；
    \item 结构化文本：包含日期、日志等级、编号等，便于正则分组替换；
    \item 边界文本：空行、仅符号行、文件末尾无换行等特殊情况。
\end{itemize}

从输入输出角度，单行处理模型可表示为：
$$L_i \xrightarrow{\text{匹配引擎}} R_i \xrightarrow{\text{替换策略}} L'_i$$
其中 $L_i$ 为第 $i$ 行原文本，$R_i$ 为命中区间集合，$L'_i$ 为替换后结果。该模型便于统一实现查找、计数、仅输出命中片段和替换写回等不同功能。

\subsection{匹配与替换算法}

本项目采用 ``双引擎匹配 + 分层替换'' 的策略。

\subsubsection{字面量匹配：KMP}

字面量模式下采用 KMP 算法。通过预处理模式串构建 LPS 数组，主串扫描时避免回退，最坏复杂度为 $\mathcal{O}\left(n+m\right)$。相较朴素匹配，KMP 在长文本、多次命中场景下更稳定。

\subsubsection{正则匹配：\mintinline{cpp}{std::regex}}

复杂模式下采用 C++23 STL 正则引擎，支持 ECMAScript 与 Extended 语法，以及捕获组反向引用替换。该方案表达能力强，适用于 ``格式化替换'' ``批量规范化'' 等任务。

\subsubsection{替换实现策略}

在字面量模式中，先得到所有命中区间，再按 ``游标前进 + 分段拼接'' 构造结果字符串，避免边替换边搜索导致的偏移错位；在正则模式中复用 \mintinline{cpp}{std::regex_replace} 语义，统一由标准库处理转义与分组替换细节。

\subsection{文件读取与写入}

\subsubsection{流式读取}

程序按行读取输入流（文件或标准输入），避免将全文一次性载入内存。该设计有两个优势：

\begin{enumerate}
    \item 对大文件更友好，内存占用稳定；
    \item 便于直接输出带行号与高亮的逐行结果。
\end{enumerate}

\subsubsection{安全写回}

原地替换时采用 ``临时文件 + 原子重命名'' 两阶段提交：先完整写入 \mintinline{text}{target.tmp}，确认写入状态后再执行 \mintinline{cpp}{std::filesystem::rename}。若失败则删除临时文件并返回错误。该策略可以避免直接覆盖带来的 ``部分写入损坏''。

\section{程序结构}

\subsection{主函数流程}

项目采用 ``薄入口、厚应用'' 的组织方式：\mintinline{text}{main.cpp} 仅调用 \mintinline{cpp}{app::run(argc, argv)}。主流程集中在应用层，整体步骤如下：

\begin{enumerate}
    \item 解析命令行参数并做约束校验；
    \item 根据参数选择 KMP 或 Regex 引擎；
    \item 初始化输入输出通道（\mintinline{text}{stdin} 或文件）；
    \item 按行执行匹配，按选项决定显示、计数或替换；
    \item 若启用原地模式，执行原子写回；
    \item 返回统一退出码（0 / 1 / 2）用于脚本集成。
\end{enumerate}

这种流程划分让 ``参数处理'' ``算法执行'' ``输出渲染'' ``写回提交'' 彼此解耦，降低维护成本。

\subsection{辅助函数设计}

辅助函数按职责分成四组：

\begin{itemize}
    \item 参数与模式辅助：完成参数互斥、依赖、位置参数归一化；
    \item 匹配辅助：提供 \mintinline{cpp}{find_all_kmp_ranges}、\mintinline{cpp}{find_all_regex_ranges} 等统一接口；
    \item 渲染辅助：负责行号拼接、命中高亮、差异展示，避免污染主循环；
    \item IO 辅助：处理输入打开、临时文件写入、原子替换与错误回滚。
\end{itemize}

该拆分策略的核心价值是 ``主流程可读、辅助函数可复用''。当后续新增目录递归或批量文件时，通常只需扩展 IO 与调度层，而无需重写匹配引擎。

\section{测试结果}

\subsection{测试环境}

构建命令如下：

\begin{minted}{bash}
cd finder
cmake -S . -B build
cmake --build build -j
\end{minted}

\subsection{测试数据}

为覆盖题目要求及常见边界，设计如下样例 \mintinline{text}{res/log.txt}：

\inputminted{text}{../res/log.txt}

并额外准备无命中文本、空文件、末尾无换行文本，用于验证退出码与写回稳定性。

\subsection{测试过程}

\subsubsection{基础字面量替换测试}

\begin{minted}{bash}
./build/finder "error" -i -r "warning" res/log.txt > res/log_kmp.txt
\end{minted}

\inputminted{text}{../res/log_kmp.txt}

\mintinline{text}{Error} / \mintinline{text}{error} 均被替换为 \mintinline{text}{warning}；其他行保持不变；输出写入 B 文件。

\subsubsection{正则分组替换测试}

\begin{minted}{bash}
./build/finder -e "(INFO|WARN|ERROR):" -i -r "[LOG:\$1]" res/log.txt > res/log_regex.txt
\end{minted}

\inputminted{text}{../res/log_regex.txt}

日志级别前缀按分组替换，分组内容保留。

\subsubsection{仅匹配输出与计数测试}

\begin{minted}{bash}
./build/finder -o -e "[0-9]{4}-[0-9]{2}-[0-9]{2}" res/log.txt
./build/finder -c "error" -i res/log.txt
\end{minted}

\inputminted{text}{../res/log_date.txt}

\inputminted{text}{../res/log_count.txt}

\mintinline{text}{-o} 仅输出匹配片段，\mintinline{text}{-c} 仅输出命中行数量。

\subsubsection{原地写回稳定性测试}

\begin{minted}{bash}
cp res/log.txt res/log_replaced.txt
./build/finder "user_" -r "account_" -p res/log_replaced.txt -d
\end{minted}

\inputminted{text}{../res/log_replaced.txt}

\mintinline{text}{log_replaced.txt} 被完整替换，未出现截断、乱码或半写入状态。

\subsection{测试结论}

综合各组测试可得：
\begin{enumerate}
    \item 程序完整满足 ``读取 A、替换、输出 B'' 的题目目标；
    \item 字面量模式在效率和稳定性上表现良好，正则模式在表达能力上优势明显；
    \item 参数互斥与依赖关系有效降低误用概率；
    \item 原地写回策略在异常场景下具备较好的数据安全性。
\end{enumerate}

\section{收获与体会}

\subsection{技术收获}

本次实验让我更清晰地理解了 ``算法正确'' 与 ``工程可用'' 之间的差异。KMP 与 Regex 解决的是匹配能力问题，而参数契约、错误码语义、原子写回解决的是工具可交付问题。只有两者结合，程序才能在真实使用中稳定运行。

\subsection{问题与解决}

开发过程中主要遇到三类问题：

\begin{itemize}
    \item 参数歧义：通过 CLI11 的互斥与依赖约束提前拦截；
    \item 输出可读性不足：引入按需颜色和 only-matching 逐命中输出；
    \item 写回风险：采用临时文件加原子替换，并在失败时回滚。
\end{itemize}

这些改进使程序从 ``能完成任务'' 提升为 ``行为可预期、结果可验证''。

\subsection{未来展望}

后续可继续从三个方向完善：

\begin{itemize}
    \item 增加目录递归与多文件批处理能力；
    \item 增加更完整的 Unicode 大小写折叠与编码兼容支持；
    \item 增加自动化回归测试与性能基准，形成持续验证流程。
\end{itemize}

\appendix

\section{程序源码}

\subsection{帮助文档}

\subsubsection{\mintinline{text}{./build/finder --help}}
\inputminted{text}{../res/help.txt}

\subsection{构建系统}

\subsubsection{\mintinline{text}{CMakeLists.txt}}
\inputminted{cmake}{../CMakeLists.txt}

\subsection{\mintinline{text}{src}}

\subsubsection{\mintinline{text}{src/main.cpp}}
\inputminted{cpp}{../src/main.cpp}

\subsubsection{\mintinline{text}{src/app.cpp}}
\inputminted{cpp}{../src/app.cpp}

\subsubsection{\mintinline{text}{src/cli.cpp}}
\inputminted{cpp}{../src/cli.cpp}

\subsubsection{\mintinline{text}{src/kmp.cpp}}
\inputminted{cpp}{../src/kmp.cpp}

\subsubsection{\mintinline{text}{src/regex.cpp}}
\inputminted{cpp}{../src/regex.cpp}

\subsection{\mintinline{text}{include}}

\subsubsection{\mintinline{text}{include/app.hpp}}
\inputminted{cpp}{../include/app.hpp}

\subsubsection{\mintinline{text}{include/cli.hpp}}
\inputminted{cpp}{../include/cli.hpp}

\subsubsection{\mintinline{text}{include/kmp.hpp}}
\inputminted{cpp}{../include/kmp.hpp}

\subsubsection{\mintinline{text}{include/regex.hpp}}
\inputminted{cpp}{../include/regex.hpp}

\subsubsection{\mintinline{text}{include/common_types.hpp}}
\inputminted{cpp}{../include/common_types.hpp}

\end{document}
